\section{本章小结}
\label{sec:matrix-approximation-summary}

本章开头的问题是：给定矩阵$\bm A$，在秩约束或正交约束下，哪一个矩阵离它最近，最小误差是多少，最优矩阵是否唯一。各节的解法遵循同一条路线。先写清约束与误差；再固定一个表示，在其中选出最优系数；最后选择表示本身，即方向、子空间或正交因子。前两步把矩阵优化化成向量的投影与特征值比较，SVD则把最优解写成具体的矩阵。

秩一逼近中，固定单位方向$\bm u$时，最优系数是投影系数$\bm c=\bm A^\top\bm u$，误差等于$\lVert\bm A\rVert_{\mathrm F}^2$减去保留量$\bm u^\top\bm A\bm A^\top\bm u$，选方向就是最大化Rayleigh商。对任意实对称$\bm S$，Rayleigh商介于最小、最大特征值之间，取等的方向恰是相应特征空间中的非零向量。用到$\bm A\bm A^\top$上，最优秩一近似为$\sigma_1\bm u_1\bm v_1^\top$，平方误差为$\sum_{i\geq2}\sigma_i^2$。它在$\sigma_1>\sigma_2$时唯一，奇异向量同时变号只改变写法。

秩$k$逼近中，固定子空间的最优重建是正交投影，选子空间化为迹极值：$k$维子空间上的迹不超过前$k$个特征值之和。由此得到Eckart--Young--Mirsky定理：截断SVD$\bm A_k$在全部秩不超过$k$的矩阵中同时使Frobenius误差与谱范数误差最小，最小值分别为$(\sum_{i>k}\sigma_i^2)^{1/2}$与$\sigma_{k+1}$。Frobenius最优解在$\sigma_k>\sigma_{k+1}$时唯一；谱范数最优解即使有间隙也一般不唯一。

这些结论只对全部秩不超过$k$的矩阵与这两种范数成立；额外的非负、稀疏约束或一般的逐元素权重，都会使截断SVD失去最优性。

改用加权长度规定单位方向时，Rayleigh商推广为广义Rayleigh商。$\bm B$正定时，用Gram因子换元可化为普通Rayleigh商，极值是广义特征值，广义特征向量可以选成$\bm B$正交。计算用Cholesky分解与三角求解保持对称结构，不应改求一般不对称的$\bm B^{-1}\bm S$。$\bm B$不正定时这一化归失效。

误差的权重能分离为行权重与列权重时，加权低秩逼近经可逆变换化为普通低秩逼近。最优解由变换后矩阵的截断SVD还原，最佳方向正是一个广义Rayleigh商的最大特征向量。这一化归要求逐元素权重是行因子与列因子的乘积；任意权重与缺失元素没有截断SVD式的闭式解，交由卷五的优化方法处理。

正交约束下，距离最小化成迹最大化，目标关于未知矩阵是线性的。换到奇异向量坐标后，目标只用到正交矩阵的对角元，最优正交矩阵是$\bm U\bm V^\top$，即第~\ref{chap:polar-decomposition}章极分解的正交因子，平方误差为$\sum_i(\sigma_i-1)^2$。$\bm A$可逆时它唯一，奇异时自由度只在零奇异值方向上。

要求旋转时，只需在最小奇异方向上作符号修正：
\[
 \bm R_\star=\bm U\operatorname{diag}\bigl(1,\ldots,1,\det(\bm U\bm V^\top)\bigr)\bm V^\top .
\]
$\det(\bm U\bm V^\top)=-1$时误差平方增加$4\sigma_d$，$\sigma_d=0$时修正不增加误差。

在这些问题中，Rayleigh商、迹极值与截断SVD承担的是最优性工具的职责：它们说明为什么由SVD读出的矩阵最好，以及在什么条件下唯一。计算出的结果要按同一种误差范数检验约束与残差，把约束带来的逼近误差与计算误差分开；最优误差的数值对扰动稳定，最优矩阵在间隙小时则不稳定。

本章的截断矩阵用于重建$\bm A$本身，与第~\ref{chap:solution-error-regularization}章为稳定求解方程而截断取逆的用途不同。不能化归到SVD的一般逼近问题，则由卷五的迭代优化展开。
